The author believes that Artificial Intelligence (AI) has been a useful tool for learning new skills,
debugging and reducing work load throughout this project.
At the start of this project, the author did not know how to program in C++.
AI systems, in particular Copilot and Claude, were used to learn about the language's syntax and features.
Moreover, the systems were helpful in learning how to use the SCIP framework.
The codebase was then designed and written by hand.
These AI systems were also used to identify errors and bugs, as well as write some utility functions.
On the data analysis side, the systems proved capable of writing and tweaking code for making plots and tables,
ensuring that these plots look as intended.
Overall, the creation of visuals was strongly aided by these systems.
The AI systems also contributed, to an extent, to finding relevant literature for \autoref{sec:ramp-pricing-problem}.
This report was written entirely by hand.
AI systems played a minor role in finding typing errors and inconsistencies, which were then resolved by hand.
In a later stage of the writing this report, AI did contribute to the proof of \autoref{thm:sep-pos}.
After the author discarded an old ``proof'' and struggled to complete a new proof based on a minimal counter-example,
AI suggested a new, relatively easy approach to formalizing the original proof.
AI also suggested ways to make the proof of \autoref{thm:sep-neg-opt} more compact.